% ============================================================
\chapter{Sous-vari\'et\'es et Seconde Forme Fondamentale}
\label{ch:sous_varietes}
% ============================================================

Comment une surface courbe se place-t-elle dans l'espace ambiant~? Cette question, au c\oe ur de la g\'eom\'etrie diff\'erentielle classique depuis Gauss, trouve sa r\'eponse moderne dans la th\'eorie des \emph{sous-vari\'et\'es riemanniennes}. La \emph{seconde forme fondamentale}, introduite par Gauss pour les surfaces dans $\R^3$ et g\'en\'eralis\'ee par Weingarten, mesure la fa\c{c}on dont la sous-vari\'et\'e ``s'incurve'' par rapport \`a l'espace ambiant. Les \'equations de Gauss, Codazzi et Ricci relient la courbure intrins\`eque de la sous-vari\'et\'e \`a sa courbure extrins\`eque. Le \emph{theorema egregium} de Gauss --- la courbure de Gauss est intrins\`eque --- en est le cas le plus c\'el\`ebre.

\section{Sous-vari\'et\'es riemanniennes}

\begin{definition}[Sous-vari\'et\'e riemannienne]
\index{sous-vari\'et\'e}
Soit $(M, g)$ une vari\'et\'e riemannienne de dimension $n$ et $\Sigma \subset M$
une sous-vari\'et\'e de dimension $k < n$. La \emph{m\'etrique induite} sur $\Sigma$
est la restriction $\bar{g} = \iota^* g$, o\`u $\iota : \Sigma \hookrightarrow M$
est l'inclusion. Le couple $(\Sigma, \bar{g})$ est appel\'e \emph{sous-vari\'et\'e
riemannienne}.
\end{definition}

\begin{definition}[D\'ecomposition du fibr\'e tangent]
\index{fibr\'e normal}
En chaque point $p \in \Sigma$, le fibr\'e tangent de $M$ se d\'ecompose orthogonalement~:
\[
    T_pM = T_p\Sigma \oplus N_p\Sigma,
\]
o\`u $N_p\Sigma = (T_p\Sigma)^\perp$ est le \emph{fibr\'e normal} de $\Sigma$ dans $M$.
On note $(\cdot)^\top$ et $(\cdot)^\perp$ les projections correspondantes.
\end{definition}

\begin{intuition}
Pensez \`a une surface $\Sigma^2$ dans $\R^3$. En chaque point, l'espace tangent \`a $\R^3$
se d\'ecompose en le plan tangent \`a $\Sigma$ et la direction normale $\nu$. La seconde
forme fondamentale mesure comment $\Sigma$ \og{}se courbe\fg dans $\R^3$, c'est-\`a-dire
comment la normale $\nu$ varie le long de $\Sigma$.
\end{intuition}

\section{Connexion induite et seconde forme fondamentale}

\begin{definition}[Formule de Gauss]
\index{formule de Gauss}
Soit $\nabla$ la connexion de Levi-Civita de $(M, g)$ et $\bar{\nabla}$ celle de
$(\Sigma, \bar{g})$. Pour tous champs de vecteurs $X, Y$ tangents \`a $\Sigma$~:
\[
    \nabla_X Y = \bar{\nabla}_X Y + \mathrm{I\!I}(X, Y),
\]
o\`u $\bar{\nabla}_X Y = (\nabla_X Y)^\top$ et $\mathrm{I\!I}(X, Y) = (\nabla_X Y)^\perp$.
\end{definition}

\begin{definition}[Seconde forme fondamentale]
\index{seconde forme fondamentale}
La \emph{seconde forme fondamentale} est l'application bilin\'eaire sym\'etrique~:
\[
    \mathrm{I\!I} : \mathfrak{X}(\Sigma) \times \mathfrak{X}(\Sigma) \to \Gamma(N\Sigma),
    \quad \mathrm{I\!I}(X, Y) = (\nabla_X Y)^\perp.
\]
\end{definition}

\begin{proposition}[Sym\'etrie de $\mathrm{I\!I}$]
La seconde forme fondamentale est sym\'etrique~: $\mathrm{I\!I}(X, Y) = \mathrm{I\!I}(Y, X)$.
\end{proposition}

\begin{proof}
On a $\mathrm{I\!I}(X, Y) - \mathrm{I\!I}(Y, X) = (\nabla_X Y - \nabla_Y X)^\perp
= [X, Y]^\perp$. Comme $X$ et $Y$ sont tangents \`a $\Sigma$, le crochet de Lie
$[X, Y]$ est aussi tangent \`a $\Sigma$, donc $[X, Y]^\perp = 0$.
\end{proof}

\begin{definition}[Op\'erateur de forme (Weingarten)]
\index{op\'erateur de Weingarten}
Pour un champ normal $\xi \in \Gamma(N\Sigma)$, l'\emph{op\'erateur de forme}
est l'endomorphisme $A_\xi : T\Sigma \to T\Sigma$ d\'efini par~:
\[
    A_\xi(X) = -(\nabla_X \xi)^\top.
\]
Il est reli\'e \`a $\mathrm{I\!I}$ par~:
$g(\mathrm{I\!I}(X, Y), \xi) = g(A_\xi(X), Y)$.
\end{definition}

\begin{proposition}[Auto-adjonction de $A_\xi$]
L'op\'erateur $A_\xi$ est auto-adjoint~:
$\bar{g}(A_\xi(X), Y) = \bar{g}(X, A_\xi(Y))$.
\end{proposition}

\begin{proof}
$\bar{g}(A_\xi(X), Y) = g(\mathrm{I\!I}(X, Y), \xi) = g(\mathrm{I\!I}(Y, X), \xi)
= \bar{g}(A_\xi(Y), X)$, par sym\'etrie de $\mathrm{I\!I}$.
\end{proof}

\section{Courbure moyenne et sous-vari\'et\'es minimales}

\begin{definition}[Courbure moyenne]
\index{courbure moyenne}
Soit $(e_1, \ldots, e_k)$ une base orthonorm\'ee de $T_p\Sigma$. Le \emph{vecteur
de courbure moyenne} est~:
\[
    \vec{H} = \frac{1}{k} \sum_{i=1}^{k} \mathrm{I\!I}(e_i, e_i)
    = \frac{1}{k} \tr(\mathrm{I\!I}).
\]
\end{definition}

\begin{definition}[Sous-vari\'et\'e minimale]
\index{sous-vari\'et\'e minimale}
Une sous-vari\'et\'e $\Sigma$ est dite \emph{minimale} si $\vec{H} = 0$, c'est-\`a-dire
si $\tr(\mathrm{I\!I}) = 0$.
\end{definition}

\begin{remarque}
Le terme \og{}minimal\fg est historique. Les sous-vari\'et\'es minimales sont les
points critiques de la fonctionnelle de volume~: $\delta \operatorname{Vol}(\Sigma) = 0$.
Elles ne minimisent pas n\'ecessairement le volume.
\end{remarque}

\begin{exemple}[Cat\'eno\"ide]
\index{cat\'eno\"ide}
La surface de r\'evolution dans $\R^3$ param\'etr\'ee par~:
\[
    \Phi(u, v) = (\cosh u \cos v, \cosh u \sin v, u), \quad (u, v) \in \R \times [0, 2\pi),
\]
est une surface minimale. On v\'erifie que $H = 0$.
\end{exemple}

\begin{exemple}[Sph\`ere $S^{n-1} \subset \R^n$]
La sph\`ere $S^{n-1}(r)$ de rayon $r$ a pour seconde forme fondamentale
$\mathrm{I\!I}(X, Y) = -\frac{1}{r} g(X, Y)\, \nu$, o\`u $\nu$ est la normale
unitaire sortante. La courbure moyenne est $\vec{H} = -\frac{1}{r}\nu$~:
la sph\`ere n'est pas minimale.
\end{exemple}

\section{\'Equations de Gauss et Codazzi}

\begin{theoreme}[\'Equation de Gauss]
\label{thm:gauss}
\index{\'equation de Gauss}
Soient $X, Y, Z, W$ des champs tangents \`a $\Sigma$. La courbure de $\Sigma$ est
reli\'ee \`a celle de $M$ par~:
\begin{multline*}
    \bar{g}(\bar{R}(X, Y)Z, W) = g(R(X, Y)Z, W) \\
    + g(\mathrm{I\!I}(X, Z), \mathrm{I\!I}(Y, W))
    - g(\mathrm{I\!I}(Y, Z), \mathrm{I\!I}(X, W)).
\end{multline*}
\end{theoreme}

\begin{proof}
Par d\'efinition, $\bar{R}(X,Y)Z = \bar{\nabla}_X \bar{\nabla}_Y Z
- \bar{\nabla}_Y \bar{\nabla}_X Z - \bar{\nabla}_{[X,Y]} Z$. En utilisant la
formule de Gauss $\nabla_X Y = \bar{\nabla}_X Y + \mathrm{I\!I}(X, Y)$ et en
d\'eveloppant~:
\begin{align*}
    \nabla_X(\bar{\nabla}_Y Z + \mathrm{I\!I}(Y,Z))
    &= \bar{\nabla}_X \bar{\nabla}_Y Z + \mathrm{I\!I}(X, \bar{\nabla}_Y Z) \\
    &\quad - A_{\mathrm{I\!I}(Y,Z)} X + \nabla_X^\perp \mathrm{I\!I}(Y,Z).
\end{align*}
En antisym\'etrisant en $X, Y$ et en projetant sur la composante tangentielle,
on obtient apr\`es simplification~:
\[
    (\bar{R}(X,Y)Z)^\top = (R(X,Y)Z)^\top + A_{\mathrm{I\!I}(X,Z)}Y - A_{\mathrm{I\!I}(Y,Z)}X.
\]
En prenant le produit scalaire avec $W \in T\Sigma$ et en utilisant
$g(A_\xi Y, W) = g(\mathrm{I\!I}(Y,W), \xi)$, on obtient le r\'esultat.
\end{proof}

\begin{corollaire}[Theorema Egregium de Gauss]
\index{Theorema Egregium}
Pour une surface $\Sigma^2 \subset \R^3$, la courbure de Gauss $K$ est intrins\`eque~:
\[
    K = \kappa_1 \kappa_2 = \det(A_\nu),
\]
o\`u $\kappa_1, \kappa_2$ sont les courbures principales.
\end{corollaire}

\begin{theoreme}[\'Equation de Codazzi]
\label{thm:codazzi}
\index{\'equation de Codazzi}
Pour tous champs $X, Y, Z$ tangents \`a $\Sigma$ et $\xi$ normal~:
\[
    g(R(X, Y)Z, \xi) = g\bigl((\bar{\nabla}_X \mathrm{I\!I})(Y, Z)
    - (\bar{\nabla}_Y \mathrm{I\!I})(X, Z),\, \xi\bigr),
\]
o\`u $(\bar{\nabla}_X \mathrm{I\!I})(Y, Z) = \nabla_X^\perp(\mathrm{I\!I}(Y, Z))
- \mathrm{I\!I}(\bar{\nabla}_X Y, Z) - \mathrm{I\!I}(Y, \bar{\nabla}_X Z)$.
\end{theoreme}

\begin{remarque}
Lorsque $M = \R^n$ (courbure nulle), l'\'equation de Codazzi se simplifie~:
$(\bar{\nabla}_X \mathrm{I\!I})(Y,Z) = (\bar{\nabla}_Y \mathrm{I\!I})(X,Z)$.
\end{remarque}

\section{Hypersurfaces}

\begin{definition}[Hypersurface]
\index{hypersurface}
Une \emph{hypersurface} est une sous-vari\'et\'e de codimension $1$. Le fibr\'e normal
est de rang $1$, engendr\'e par un champ unitaire $\nu$.
\end{definition}

Pour une hypersurface, la seconde forme fondamentale scalaire s'\'ecrit~:
\[
    h(X, Y) = g(\mathrm{I\!I}(X, Y), \nu) = g(A_\nu(X), Y).
\]

\begin{definition}[Courbures principales]
\index{courbures principales}
Les \emph{courbures principales} $\kappa_1, \ldots, \kappa_{n-1}$ sont les valeurs
propres de $A_\nu$. Les vecteurs propres sont les \emph{directions principales}.
\end{definition}

\begin{proposition}
Pour une hypersurface~:
\begin{align*}
    H &= \frac{1}{n-1}(\kappa_1 + \cdots + \kappa_{n-1})
       = \frac{1}{n-1} \tr(A_\nu), \\
    K_{\text{Gauss}} &= \kappa_1 \cdots \kappa_{n-1} = \det(A_\nu).
\end{align*}
\end{proposition}

\section{Sous-vari\'et\'es totalement g\'eod\'esiques et ombilicales}

\begin{definition}[Totalement g\'eod\'esique]
\index{totalement g\'eod\'esique}
$\Sigma$ est \emph{totalement g\'eod\'esique} si $\mathrm{I\!I} \equiv 0$.
\end{definition}

\begin{exemple}
Les grands cercles de $S^2$ et les \'equateurs $S^{k} \subset S^n$ sont totalement
g\'eod\'esiques.
\end{exemple}

\begin{definition}[Totalement ombilicale]
\index{totalement ombilicale}
$\Sigma$ est \emph{totalement ombilicale} si $\mathrm{I\!I}(X, Y) = g(X, Y)\, \vec{H}$.
\end{definition}

\begin{exemple}
Les sph\`eres $S^{n-1}(r) \subset \R^n$ sont totalement ombilicales, avec
$\kappa_1 = \cdots = \kappa_{n-1} = 1/r$.
\end{exemple}

\begin{formules}[title=\textbf{Formules fondamentales des sous-vari\'et\'es}]
\begin{align*}
    \text{Gauss~:} \quad & \nabla_X Y = \bar{\nabla}_X Y + \mathrm{I\!I}(X, Y) \\
    \text{Weingarten~:} \quad & \nabla_X \xi = -A_\xi(X) + \nabla_X^\perp \xi \\
    \text{\'Eq. de Gauss~:} \quad & \bar{g}(\bar{R}(X,Y)Z, W) = g(R(X,Y)Z, W)
        + g(\mathrm{I\!I}(X,Z), \mathrm{I\!I}(Y,W)) \\
    & \qquad\qquad\qquad\quad - g(\mathrm{I\!I}(Y,Z), \mathrm{I\!I}(X,W)) \\
    \text{Courbure moy.~:} \quad & \vec{H} = \frac{1}{k}\tr(\mathrm{I\!I})
\end{align*}
\end{formules}

\section{Exercices}

\begin{exercice}
Calculer la seconde forme fondamentale du tore de r\'evolution
$\Phi(\theta, \phi) = ((R + r\cos\theta)\cos\phi,
(R + r\cos\theta)\sin\phi, r\sin\theta)$ dans $\R^3$. D\'eterminer les courbures
principales et la courbure moyenne.
\end{exercice}

\begin{exercice}
Montrer que pour une hypersurface $\Sigma^{n-1} \subset \R^n$, l'\'equation de Gauss
donne $\bar{K}(\sigma) = \kappa_i \kappa_j$ pour le plan $\sigma$ engendr\'e par les
directions principales $e_i$ et $e_j$.
\end{exercice}

\begin{exercice}
Soit $\Sigma$ une surface minimale dans $\R^3$. Montrer que $\kappa_1 = -\kappa_2$
en tout point, et que $K = -\kappa_1^2 \leq 0$.
\end{exercice}

\begin{exercice}
Montrer que l'h\'elico\"ide $\Phi(u,v) = (v\cos u, v\sin u, au)$ est une surface
minimale dans $\R^3$.
\end{exercice}

\begin{exercice}
Soit $\Sigma$ une hypersurface compacte sans bord de $\R^n$. Montrer qu'il existe
un point $p \in \Sigma$ o\`u toutes les courbures principales sont strictement positives.
\emph{Indication~:} consid\'erer le point le plus \'eloign\'e de l'origine.
\end{exercice}

\begin{exercice}
Soit $f : \R^2 \to \R$ une fonction lisse et $\Sigma = \{(x, y, f(x,y))\}$ le graphe
de $f$ dans $\R^3$. Calculer la seconde forme fondamentale, les courbures principales,
et la courbure moyenne de $\Sigma$ en fonction de $f$ et ses d\'eriv\'ees.
\end{exercice}
